%% =============================================================================
%%  Etapa 4 — Contexto
%% =============================================================================
\chapter{Contexto}\label{cap:contexto}

\begin{objetivo}
Situar o trecho em três círculos: o discurso de que faz parte, o Evangelho
de Lucas como um todo e o paralelo em Mateus. Por fim, perguntar pelo mundo
social em que \enquote{pobre}, \enquote{rico}, \enquote{honra} e
\enquote{profeta} tinham sentido.
\end{objetivo}

\section{Contexto imediato}

\begin{ebtabela}{@{}l >{\raggedright\arraybackslash}X@{}}
\EBcab{Trecho} & \EBcab{Conteúdo}\\\midrule
\rb{6:12–16} & Jesus passa a noite em oração no monte e escolhe os Doze\\
\rb{6:17–19} & desce a um \enquote{lugar plano}; multidão de discípulos e de povo da Judeia, Jerusalém, Tiro e Sidom; curas\\
\rb{6:20–26} & \textbf{felicidades e ais}, dirigidos aos discípulos (\enquote{erguendo os olhos para os seus discípulos})\\
\rb{6:27–38} & amor aos inimigos, regra de ouro, não julgar\\
\rb{6:39–45} & parábolas: cego guiando cego, árvore e fruto\\
\rb{6:46–49} & ouvir e praticar: as duas casas\\
\end{ebtabela}

Dois dados orientam a leitura. Os destinatários são os \emph{discípulos}
(v.~20), mas a multidão ouve (\rb{6:17–19; 7:1}). E o discurso começa por
declarar \emph{quem Deus favorece} antes de dizer \emph{o que fazer}
(vv.~27ss.): a ética vem depois da boa notícia.

\section{Temas lucanos}

O trecho condensa temas que atravessam todo o Evangelho:

\begin{ebtabela}{@{}l >{\raggedright\arraybackslash}X@{}}
\EBcab{Texto} & \EBcab{Ligação com Lc 6:20–26}\\\midrule
\rb{1:51–53} & Magnificat: derruba poderosos, eleva humildes; enche de bens os famintos, despede vazios os ricos (\gr{ἐνέπλησεν})\\
\rb{2:25} & Simeão espera a \gr{παράκλησις} de Israel\\
\rb{4:18–19} & programa de Nazaré: \enquote{boa notícia aos pobres} (\rb{Is 61:1–2})\\
\rb{7:22} & resposta a João Batista: \enquote{aos pobres é anunciado o evangelho}\\
\rb{12:13–21} & o rico insensato: \enquote{descansa, come, bebe e alegra-te}\\
\rb{14:12–24} & convidar pobres, aleijados, coxos e cegos; o grande banquete\\
\rb{16:19–31} & o rico e Lázaro: \enquote{recebeste os teus bens em tua vida\ldots\ agora ele é consolado} (16:25)\\
\rb{18:18–30} & o jovem rico; \enquote{como é difícil aos que têm riquezas\ldots}\\
\rb{19:1–10} & Zaqueu dá metade dos bens aos pobres\\
\end{ebtabela}

\begin{pergunta}
Compare \rb{6:24} com \rb{16:25}. Que palavras e ideias se repetem? O que a
parábola acrescenta ao \enquote{ai} dirigido aos ricos?
\end{pergunta}

\section{Leitura sinótica}

Mateus traz nove macarismos (\rb{Mt 5:3–12}) e nenhum ai; Lucas traz quatro
macarismos e quatro ais. A comparação não serve para \enquote{harmonizar} os
dois, mas para perceber a ênfase de cada um.

\begingroup\small
\begin{ebtabela}{@{}>{\raggedright}p{.46\linewidth} >{\raggedright\arraybackslash}X@{}}
\EBcab{Lucas 6} & \EBcab{Mateus 5}\\\midrule
\gr{Μακάριοι οἱ πτωχοί} — 2.ª pessoa (\enquote{vosso é o reino}) &
  \gr{Μακάριοι οἱ πτωχοὶ τῷ πνεύματι} — 3.ª pessoa (\enquote{deles é o reino dos céus})\\
\gr{οἱ πεινῶντες νῦν} & \gr{οἱ πεινῶντες καὶ διψῶντες τὴν δικαιοσύνην}\\
\gr{οἱ κλαίοντες νῦν, ὅτι γελάσετε} & \gr{οἱ πενθοῦντες, ὅτι αὐτοὶ παρακληθήσονται}\\
\gr{μισήσωσιν, ἀφορίσωσιν, ὀνειδίσωσιν, ἐκβάλωσιν} & \gr{ὀνειδίσωσιν, διώξωσιν, εἴπωσιν πᾶν πονηρόν}\\
\gr{ἕνεκα τοῦ υἱοῦ τοῦ ἀνθρώπου} & \gr{ἕνεκεν ἐμοῦ}\\
\gr{χάρητε \ldots\ καὶ σκιρτήσατε} & \gr{χαίρετε καὶ ἀγαλλιᾶσθε}\\
\gr{ἐν τῷ οὐρανῷ} & \gr{ἐν τοῖς οὐρανοῖς}\\
\gr{ἐποίουν} (faziam) & \gr{ἐδίωξαν} (perseguiram)\\
quatro ais (vv.~24–26) & — (ais contra escribas e fariseus só em \rb{Mt 23})\\
\end{ebtabela}
\endgroup

\paragraph{O que a tabela mostra.} Em Lucas, a pobreza, a fome e o choro
aparecem sem qualificação e no \enquote{agora}; em Mateus, com dimensão
interior e ética (\enquote{em espírito}, \enquote{de justiça}). É comum ler
isso como oposição (Lucas \enquote{social}, Mateus \enquote{espiritual}), mas
em ambos os destinatários são discípulos e em ambos a promessa depende de
Deus. Note ainda que o \gr{παρακληθήσονται} de Mateus ecoa a
\gr{παράκλησις} que os ricos \enquote{já receberam} em Lucas.

\begin{alerta}[Harmonização]
Não é preciso decidir qual versão é \enquote{a original} para entender cada
uma. Primeiro leia Lucas como Lucas; o diálogo com Mateus vem depois
\parencite{fee2002}.
\end{alerta}

\section{Contexto social e histórico}

\paragraph{Honra e vergonha.} No mundo mediterrâneo antigo, a honra
(reconhecimento público) valia tanto quanto os bens. \textcite{hanson1994}
lê os macarismos como declarações de honra e os ais como declarações de
vergonha. Os vv.~22 e~26 falam justamente disso: ser excluído e ter o nome
rejeitado (vergonha pública) ou ser elogiado por todos (honra pública).

\paragraph{Exclusão.} \gr{ἀφορίζω} (\enquote{separar}, \enquote{excluir})
pode aludir à exclusão da sinagoga, experiência conhecida das primeiras
comunidades (\rb{Jo 9:22; 16:2}).

\paragraph{Profetas perseguidos, falsos profetas aplaudidos.} A tradição de
Israel lembrava que os profetas eram rejeitados (\rb{Ne 9:26; 2Cr 36:15–16};
cf.\ \rb{Lc 11:47–51; 13:33–34; At 7:52}) e que os falsos profetas diziam o
que o povo queria ouvir (\rb{Jr 5:31; 6:14; Mq 2:11}).

\begin{pergunta}\label{perg:pobres}
Pobres e ricos, em Lucas, são categorias econômicas, religiosas ou as duas?
Liste os textos da tabela de temas lucanos que apoiam cada leitura antes de
consultar os comentários.
\end{pergunta}
